% Slide: Quy Trình Hiện Tại (As-Is)
\begin{frame}{Quy Trình Hiện Tại (As-Is)}
  \framesubtitle{\color{white!90}\textbf{Hành trình nhiều bất tiện và rủi ro của chủ nuôi}}

  \begin{columns}[T]
    % Cột 1: 4 Đoạn hành trình
    \begin{column}{0.50\textwidth}
      {\color{PetcarePrimary}\large \textbf{4 Chặng quy trình thủ công}}\\[2mm]
      \small
      \begin{itemize}
        \setlength\itemsep{2.8mm}
        \item \textbf{1. Đặt lịch hẹn:} Gọi điện/chat hỏi giờ trống; nhân viên lật sổ tay đối soát; chờ phản hồi 15p--2h.
        \item \textbf{2. Dặn dò dị ứng:} Dặn miệng hoặc dán post-it lên chuồng; dễ rơi mất khi đổi ca trực.
        \item \textbf{3. Chờ đợi (Điểm mù):} Thú cưng ở tiệm 2--3 tiếng; chủ nuôi bồn chồn lo âu nhưng ngại gọi.
        \item \textbf{4. Đón bé \& Thanh toán:} Nhận cuộc gọi gián đoạn công việc; hóa đơn giấy dễ mất, quên phác đồ cũ.
      \end{itemize}
    \end{column}

    % Cột 2: Điểm nghẽn & Rủi ro chí mạng dạng Bullet
    \begin{column}{0.46\textwidth}
      {\color{PetcarePrimary}\large \textbf{4 Điểm nghẽn \& Rủi ro chí mạng}}\\[2mm]
      \small
      \begin{itemize}
        \setlength\itemsep{2.8mm}
        \item \textbf{Độ trễ đặt lịch (15 -- 120p):} Tra sổ tay thủ công, dễ bị trùng lịch hoặc lỡ mất khung giờ đẹp.
        \item \textbf{Bỏ sót dị ứng ($>$ 80\% rủi ro):} Ghi nhớ cảm tính, không có cơ chế đối soát bắt buộc khi đổi thợ.
        \item \textbf{Mù mờ tiến độ (2 -- 3 tiếng):} Đứt gãy luồng thông tin, tạo tâm lý lo âu kéo dài cho chủ nuôi.
        \item \textbf{Mất dấu lịch sử (100\%):} Hóa đơn và tên dầu tắm trôi dạt, không theo dõi được sức khỏe.
      \end{itemize}
    \end{column}
  \end{columns}
\end{frame}
